% ===== PRÉAMBULE CANONIQUE — à coller VERBATIM en tête de CHAQUE .tex =====
\documentclass[11pt,a4paper]{article}
\usepackage[utf8]{inputenc}\usepackage[T1]{fontenc}\usepackage{lmodern}
\usepackage[french]{babel}
\usepackage{amsmath,amssymb,mathtools}\usepackage{icomma}
\usepackage[version=4]{mhchem}          % \ce{...} : équations & formules
\usepackage{siunitx}\sisetup{locale=FR} % \qty{}{}, \si{}, \ang{}
\usepackage{chemfig}                    % \chemfig{...} : structures développées
\usepackage{xcolor}\usepackage{colortbl}
\usepackage{tikz}\usepackage{pgfplots}\pgfplotsset{compat=1.18}
\usepackage[most]{tcolorbox}\usepackage{enumitem}\usepackage{array,booktabs}
\usepackage[a4paper,margin=2.2cm]{geometry}
\usetikzlibrary{arrows.meta,calc,angles,quotes,positioning,patterns,backgrounds,shapes.geometric}
\definecolor{primary}{RGB}{222,49,99}\definecolor{primarydark}{RGB}{150,22,55}
\definecolor{defcol}{RGB}{0,114,178}\definecolor{methcol}{RGB}{52,92,148}
\definecolor{excol}{RGB}{0,135,90}\definecolor{attncol}{RGB}{200,40,55}
\tcbset{boxbase/.style={enhanced,breakable,boxrule=0.9pt,arc=2.5pt,
  left=3mm,right=3mm,top=2.3mm,bottom=2.3mm,fonttitle=\bfseries,coltitle=white}}
% --- boîtes TUTORIEL (noms EXACTS, ne pas renommer) ---
\newtcolorbox{cle}[1][]{boxbase,colback=defcol!6,colframe=defcol,
  title={\textsc{À retenir}\ifx\relax#1\relax\relax\else\textmd{ : \itshape #1}\fi}}
\newtcolorbox{methode}[1][]{boxbase,colback=methcol!7,colframe=methcol,sharp corners,
  title={\textsc{Méthode}\ifx\relax#1\relax\relax\else\textmd{ --- \itshape #1}\fi}}
\newtcolorbox{exemplebox}[1][]{boxbase,colback=excol!5,colframe=excol,
  title={\textsc{Exemple}\ifx\relax#1\relax\relax\else\textmd{ : \itshape #1}\fi}}
\newtcolorbox{piege}[1][]{boxbase,colback=attncol!6,colframe=attncol,
  title={\textsc{Piège}\ifx\relax#1\relax\relax\else\textmd{ : \itshape #1}\fi}}
\newtcolorbox{remarque}[1][]{boxbase,colback=black!4,colframe=black!45,
  title={\textsc{Remarque}\ifx\relax#1\relax\relax\else\textmd{ : \itshape #1}\fi}}
% --- boîtes EXERCICE / CORRIGÉ (noms EXACTS, auto-comptées) ---
\newtcolorbox[auto counter]{exo}[1][]{boxbase,colback=primary!4,colframe=primary,
  title={\textsc{Exercice}~\thetcbcounter\ifx\relax#1\relax\relax\else\textmd{ --- \itshape #1}\fi}}
\newtcolorbox[auto counter]{corrige}[1][]{boxbase,colback=excol!4,colframe=excol,
  title={\textsc{Corrigé}~\thetcbcounter\ifx\relax#1\relax\relax\else\textmd{ --- \itshape #1}\fi}}
\newcommand{\R}{\mathbb{R}}\newcommand{\N}{\mathbb{N}}\newcommand{\Z}{\mathbb{Z}}\newcommand{\ds}{\displaystyle}
% (un \usepackage{fancyhdr}+en-tête est possible ; il est ignoré côté web)
% =========================================================================
\begin{document}
\begin{center}\color{primarydark}\large\bfseries Thème 4 — La formule de Nernst \quad$\bullet$\quad Niveau Facile\end{center}

\begin{exo}[les réflexes de calcul]
Avant d'appliquer Nernst, entraîne-toi.
\begin{enumerate}[label=\alph*)]
  \item Calcule $\log(100)$, $\log(10)$, $\log(1)$, $\log(0{,}1)$ et $\log(0{,}01)$.
  \item Calcule le facteur $\dfrac{0{,}06}{n}$ pour $n=1$, $n=2$ et $n=3$.
\end{enumerate}
\end{exo}

\begin{exo}[quand le potentiel vaut $E_{\mathrm{r}}^{\circ}$]
On rappelle $E=E_{\mathrm{r}}^{\circ}+\dfrac{0{,}06}{n}\log\dfrac{[\mathrm{Ox}]}{[\mathrm{Red}]}$.
\begin{enumerate}[label=\alph*)]
  \item Couple \ce{Fe^{3+}}/\ce{Fe^{2+}} ($E_{\mathrm{r}}^{\circ}=+0{,}77$~V, $n=1$) avec
        $[\ce{Fe^{3+}}]=[\ce{Fe^{2+}}]$. Que vaut $E$ ?
  \item Couple \ce{Cu^{2+}}/\ce{Cu} ($E_{\mathrm{r}}^{\circ}=+0{,}34$~V, $n=2$) avec
        $[\ce{Cu^{2+}}]=1$~mol/L. Que vaut $E$ ?
\end{enumerate}
\end{exo}

\begin{exo}[électrode diluée]
Électrode de cuivre \ce{Cu^{2+} + 2 e^- <=> Cu} ($E_{\mathrm{r}}^{\circ}=+0{,}34$~V, $n=2$) plongée
dans une solution où $[\ce{Cu^{2+}}]=0{,}01$~mol/L. Calcule le potentiel $E$ de la demi-pile.
\end{exo}

\begin{exo}[force électromotrice standard]
Une pile fonctionne en conditions standard ($[\,]=1$~mol/L). Données :
$E_{\mathrm{r}}^{\circ}(\ce{Cu^{2+}}/\ce{Cu})=+0{,}34$~V et
$E_{\mathrm{r}}^{\circ}(\ce{Zn^{2+}}/\ce{Zn})=-0{,}76$~V.
\begin{enumerate}[label=\alph*)]
  \item Quelle demi-pile est la cathode ? l'anode ?
  \item Calcule la f.é.m.\ standard $\Delta E^{\circ}=E_{\mathrm{r}}^{\circ}(\text{cathode})
        -E_{\mathrm{r}}^{\circ}(\text{anode})$.
\end{enumerate}
\end{exo}

\begin{exo}[cathode, anode, f.é.m.]
Pile utilisant les couples \ce{Ag+}/\ce{Ag} ($+0{,}80$~V) et \ce{Fe^{2+}}/\ce{Fe} ($-0{,}44$~V), en
conditions standard.
\begin{enumerate}[label=\alph*)]
  \item Identifie la cathode et l'anode.
  \item Calcule la f.é.m.\ standard.
\end{enumerate}
\end{exo}

\end{document}
